\subsubsection{\texttt{aveknt}}

\texttt{aveknt}用于计算节点数组元素的平均值。具体来说，它根据给定的节点序列 $t$ 和样条曲线的阶数 $k$，返回连续 $k-1$ 个节点的平均值作为新的节点位置。

用法为：\texttt{tstar = aveknt( t, k )}
注意：在k = t的长度的时候，返回的新节点序列长度为tatsr为空向量。
\begin{itemize}
    \item \texttt{t} 为一个单调递增的节点序列。
    \item \texttt{k} 为样条曲线的阶数，$k \geq 2$且k为整数。
    \item \texttt{tstar} 为生成的新节点序列。
\end{itemize}
生成算法为：$t_{i}^{*} := \frac{t_{i+1} + \cdots + t_{i+k-1}}{k-1}, \quad i = 1:n$， 其中n为节点序列t的长度，$t_{i}^{*}$为新的节点序列。
输入：
\begin{lstlisting}[language = matlab]
    x = [1,2,3,3,3];
    k = 3;
    tstar = aveknt(x, k);
\end{lstlisting}
输出：
\begin{lstlisting}[language = matlab]
    tstar = 

    2.5000    3.0000
\end{lstlisting}

输入：
\begin{lstlisting}[language = matlab]
    y = [1,2,3];
    k = 3;
    tstar = aveknt(y, k);
\end{lstlisting}
输出：
\begin{lstlisting}[language = matlab]
    tstar = 

    1x0 empty double

\end{lstlisting}

